% Options for packages loaded elsewhere
\PassOptionsToPackage{unicode}{hyperref}
\PassOptionsToPackage{hyphens}{url}
%
\documentclass[
  11pt,
]{article}
\usepackage{amsmath,amssymb}
\usepackage{lmodern}
\usepackage{iftex}
\ifPDFTeX
  \usepackage[T1]{fontenc}
  \usepackage[utf8]{inputenc}
  \usepackage{textcomp} % provide euro and other symbols
\else % if luatex or xetex
  \usepackage{unicode-math}
  \defaultfontfeatures{Scale=MatchLowercase}
  \defaultfontfeatures[\rmfamily]{Ligatures=TeX,Scale=1}
  \setmainfont[]{Times New Roman}
\fi
% Use upquote if available, for straight quotes in verbatim environments
\IfFileExists{upquote.sty}{\usepackage{upquote}}{}
\IfFileExists{microtype.sty}{% use microtype if available
  \usepackage[]{microtype}
  \UseMicrotypeSet[protrusion]{basicmath} % disable protrusion for tt fonts
}{}
\makeatletter
\@ifundefined{KOMAClassName}{% if non-KOMA class
  \IfFileExists{parskip.sty}{%
    \usepackage{parskip}
  }{% else
    \setlength{\parindent}{0pt}
    \setlength{\parskip}{6pt plus 2pt minus 1pt}}
}{% if KOMA class
  \KOMAoptions{parskip=half}}
\makeatother
\usepackage{xcolor}
\usepackage[margin=1in]{geometry}
\usepackage{color}
\usepackage{fancyvrb}
\newcommand{\VerbBar}{|}
\newcommand{\VERB}{\Verb[commandchars=\\\{\}]}
\DefineVerbatimEnvironment{Highlighting}{Verbatim}{commandchars=\\\{\}}
% Add ',fontsize=\small' for more characters per line
\usepackage{framed}
\definecolor{shadecolor}{RGB}{248,248,248}
\newenvironment{Shaded}{\begin{snugshade}}{\end{snugshade}}
\newcommand{\AlertTok}[1]{\textcolor[rgb]{0.94,0.16,0.16}{#1}}
\newcommand{\AnnotationTok}[1]{\textcolor[rgb]{0.56,0.35,0.01}{\textbf{\textit{#1}}}}
\newcommand{\AttributeTok}[1]{\textcolor[rgb]{0.77,0.63,0.00}{#1}}
\newcommand{\BaseNTok}[1]{\textcolor[rgb]{0.00,0.00,0.81}{#1}}
\newcommand{\BuiltInTok}[1]{#1}
\newcommand{\CharTok}[1]{\textcolor[rgb]{0.31,0.60,0.02}{#1}}
\newcommand{\CommentTok}[1]{\textcolor[rgb]{0.56,0.35,0.01}{\textit{#1}}}
\newcommand{\CommentVarTok}[1]{\textcolor[rgb]{0.56,0.35,0.01}{\textbf{\textit{#1}}}}
\newcommand{\ConstantTok}[1]{\textcolor[rgb]{0.00,0.00,0.00}{#1}}
\newcommand{\ControlFlowTok}[1]{\textcolor[rgb]{0.13,0.29,0.53}{\textbf{#1}}}
\newcommand{\DataTypeTok}[1]{\textcolor[rgb]{0.13,0.29,0.53}{#1}}
\newcommand{\DecValTok}[1]{\textcolor[rgb]{0.00,0.00,0.81}{#1}}
\newcommand{\DocumentationTok}[1]{\textcolor[rgb]{0.56,0.35,0.01}{\textbf{\textit{#1}}}}
\newcommand{\ErrorTok}[1]{\textcolor[rgb]{0.64,0.00,0.00}{\textbf{#1}}}
\newcommand{\ExtensionTok}[1]{#1}
\newcommand{\FloatTok}[1]{\textcolor[rgb]{0.00,0.00,0.81}{#1}}
\newcommand{\FunctionTok}[1]{\textcolor[rgb]{0.00,0.00,0.00}{#1}}
\newcommand{\ImportTok}[1]{#1}
\newcommand{\InformationTok}[1]{\textcolor[rgb]{0.56,0.35,0.01}{\textbf{\textit{#1}}}}
\newcommand{\KeywordTok}[1]{\textcolor[rgb]{0.13,0.29,0.53}{\textbf{#1}}}
\newcommand{\NormalTok}[1]{#1}
\newcommand{\OperatorTok}[1]{\textcolor[rgb]{0.81,0.36,0.00}{\textbf{#1}}}
\newcommand{\OtherTok}[1]{\textcolor[rgb]{0.56,0.35,0.01}{#1}}
\newcommand{\PreprocessorTok}[1]{\textcolor[rgb]{0.56,0.35,0.01}{\textit{#1}}}
\newcommand{\RegionMarkerTok}[1]{#1}
\newcommand{\SpecialCharTok}[1]{\textcolor[rgb]{0.00,0.00,0.00}{#1}}
\newcommand{\SpecialStringTok}[1]{\textcolor[rgb]{0.31,0.60,0.02}{#1}}
\newcommand{\StringTok}[1]{\textcolor[rgb]{0.31,0.60,0.02}{#1}}
\newcommand{\VariableTok}[1]{\textcolor[rgb]{0.00,0.00,0.00}{#1}}
\newcommand{\VerbatimStringTok}[1]{\textcolor[rgb]{0.31,0.60,0.02}{#1}}
\newcommand{\WarningTok}[1]{\textcolor[rgb]{0.56,0.35,0.01}{\textbf{\textit{#1}}}}
\usepackage{longtable,booktabs,array}
\usepackage{calc} % for calculating minipage widths
% Correct order of tables after \paragraph or \subparagraph
\usepackage{etoolbox}
\makeatletter
\patchcmd\longtable{\par}{\if@noskipsec\mbox{}\fi\par}{}{}
\makeatother
% Allow footnotes in longtable head/foot
\IfFileExists{footnotehyper.sty}{\usepackage{footnotehyper}}{\usepackage{footnote}}
\makesavenoteenv{longtable}
\usepackage{graphicx}
\makeatletter
\def\maxwidth{\ifdim\Gin@nat@width>\linewidth\linewidth\else\Gin@nat@width\fi}
\def\maxheight{\ifdim\Gin@nat@height>\textheight\textheight\else\Gin@nat@height\fi}
\makeatother
% Scale images if necessary, so that they will not overflow the page
% margins by default, and it is still possible to overwrite the defaults
% using explicit options in \includegraphics[width, height, ...]{}
\setkeys{Gin}{width=\maxwidth,height=\maxheight,keepaspectratio}
% Set default figure placement to htbp
\makeatletter
\def\fps@figure{htbp}
\makeatother
\setlength{\emergencystretch}{3em} % prevent overfull lines
\providecommand{\tightlist}{%
  \setlength{\itemsep}{0pt}\setlength{\parskip}{0pt}}
\setcounter{secnumdepth}{-\maxdimen} % remove section numbering
\ifLuaTeX
  \usepackage{selnolig}  % disable illegal ligatures
\fi
\IfFileExists{bookmark.sty}{\usepackage{bookmark}}{\usepackage{hyperref}}
\IfFileExists{xurl.sty}{\usepackage{xurl}}{} % add URL line breaks if available
\urlstyle{same} % disable monospaced font for URLs
\hypersetup{
  pdftitle={Phân Cụm Cộng Đồng Game Trên Nền Tảng Steam - Dự Đoán Mối Quan Hệ Kết Bạn Giữa Game Thủ Steam},
  pdfauthor={Tên bạn},
  hidelinks,
  pdfcreator={LaTeX via pandoc}}

\title{Phân Cụm Cộng Đồng Game Trên Nền Tảng Steam - Dự Đoán Mối Quan Hệ
Kết Bạn Giữa Game Thủ Steam}
\author{Tên bạn}
\date{2025-09-29}

\begin{document}
\maketitle

\hypertarget{i.-giux1edbi-thiux1ec7u-vuxe0-thiux1ebft-lux1eadp}{%
\section{I. Giới Thiệu và Thiết
Lập}\label{i.-giux1edbi-thiux1ec7u-vuxe0-thiux1ebft-lux1eadp}}

Nghiên cứu này là một phân tích nhằm khám phá cấu trúc và động lực xã
hội của cộng đồng game thủ trên nền tảng Steam. Mục tiêu cốt lõi là định
hình thành hai giai đoạn liên kết:

Phân Cụm Cộng Đồng (Clustering): Sử dụng các đặc trưng đa chiều (thời
gian chơi game, thành tích, nhóm tham gia, và tương tác xã hội) để xác
định và định danh các cộng đồng người dùng có hành vi đồng nhất.

Dự Đoán Mối Quan Hệ (Link Prediction): Ứng dụng kết quả phân cụm cộng
đồǹg và các biến dữ liệu cá nhân, làm biến dự báo để xây dựng mô hình dự
đoán khả năng kết bạn giữa các cặp game thủ trên mạng lưới Steam.

\begin{figure}
\centering
\includegraphics{image/p1.jpg}
\caption{Mô tả}
\end{figure}

\hypertarget{ii.-thu-thux1eadp-dux1eef-liux1ec7u-data-acquisition}{%
\section{II. Thu thập Dữ liệu (Data
Acquisition)}\label{ii.-thu-thux1eadp-dux1eef-liux1ec7u-data-acquisition}}

Phần này mô tả chi tiết quy trình ** thu thập dữ liệu** từ Steam Web API
và tạo tập dữ liệu thô. Vì để vào steam communty để cào dữ liệu, người
dùng sẽ bị chặn hoặc găp nhiều khó khăn trong quá trình thu thập như bị
chặn hoặc capcha,.. .Chỉ có một cách là dùng account nào đủ uy tín ( đã
mua game hoạc trade game có giao dịch trên 5usd sẽ được cấp Steam Web
API Key). Sau khi có Key thì thực hiện thu thập dữ liệu

\hypertarget{cuxe0o-dux1eef-liux1ec7u-tux1eeb-steam-web-api}{%
\subsection{2.1 Cào dữ liệu từ Steam Web
API}\label{cuxe0o-dux1eef-liux1ec7u-tux1eeb-steam-web-api}}

Dữ liệu được thu thập thông qua \textbf{Steam Web API}, một giao diện
cho phép truy cập trực tiếp vào các thông tin công khai của người dùng
và hệ sinh thái trên nền tảng Steam.

Quá trình thu thập dữ liệu trải qua nhiều giai đoạn và thu thập nhiều
loại thông tin. Đầu tiên, vì Steam Web API không cho phép cào dữ liệu
một cách ngẫu nhiên, nên cần đưa vào steam id của một người dùng nào đó
trên steam để API tìm, khi đó chương trình sẽ dẫn đến trang cá nhân của
người đó.

Trong một trang cá nhân bao gồm nhiều thông tin như: id, tên, thời gian
tham gia, nước, khu vực, số thành tựu, số trò chơi đã có, list bạn bè và
nhóm tham gia,\ldots{} Để thu thập các thông tin trên ta sử dụng API
endpoints( bao gồm các biến mà steam cung cấp để lấy các dữ liệu tương
ứng):

\begin{itemize}
\tightlist
\item
  GetFriendList : Lấy danh sách bạn bè của một người
\item
  GetUserGroupList : Lấy danh sách nhóm của một người
\item
  get\_owned\_games : Lấy danh sách game của một người
\item
  get\_owned\_achievements: lấy danh sách thành tựu của một người
\end{itemize}

Và chỉ thu thập thông tin của người nào đặt trạng thái là public, các
account để Private thì bo qua. Sau khi thu thập xong một trang cá nhân
thì chương trình sẽ chuyển hướng đến danh sách bạn bè của người đó, thực
hiện truy cập vào trang cá nhân của người đó và tiếp tục thu
thập.~Chương trình sẽ tiếp tục thu thập theo hướng đó cho tới khi đủ
2000 người dùng public.

\begin{figure}
\centering
\includegraphics{image/p2.png}
\caption{Dữ liệu sau khi đã được crawl}
\end{figure}

\begin{figure}
\centering
\includegraphics{image/p4.png}
\caption{WorkFlow Hậu thu thập dữ liệu}
\end{figure}

\begin{longtable}[]{@{}
  >{\raggedright\arraybackslash}p{(\columnwidth - 4\tabcolsep) * \real{0.3913}}
  >{\raggedright\arraybackslash}p{(\columnwidth - 4\tabcolsep) * \real{0.3913}}
  >{\raggedright\arraybackslash}p{(\columnwidth - 4\tabcolsep) * \real{0.2174}}@{}}
\toprule()
\begin{minipage}[b]{\linewidth}\raggedright
Hạng mục
\end{minipage} & \begin{minipage}[b]{\linewidth}\raggedright
Dữ liệu thu thập
\end{minipage} & \begin{minipage}[b]{\linewidth}\raggedright
Mục đích
\end{minipage} \\
\midrule()
\endhead
Profiles\_Users & Thông tin cơ bản:
\texttt{timecreated},\texttt{user\_id},
\texttt{loccountrycode},\texttt{personaname}, \texttt{loccountrycode}. &
Tính tuổi tài khoản, vị trí địa lý. \\
Achievements & Danh sách achievements và thời gian mở khóa. & Đo lường
hành vi chơi game \\
Friends & Danh sách bạn bè (\texttt{,friend\_steamids\_list},
\texttt{friend\_count}). & Dữ liệu chính cho dự đoán mối quan hệ \\
Games & Danh sách game sở hữu, \texttt{playtime\_forever}. & Tổng quan
về dữ liệu trò chơi của người dùng \\
Groups & Danh sách nhóm tham gia
(\texttt{filtered\_all\_groups\_gid\_30},\texttt{group\_names}). & Dùng
cho phân cụm cộng đồng. \\
\bottomrule()
\end{longtable}

\hypertarget{gux1ed9p-vuxe0-chuux1ea9n-bux1ecb-dux1eef-liux1ec7u}{%
\subsection{2.2 Gộp và Chuẩn bị Dữ
liệu}\label{gux1ed9p-vuxe0-chuux1ea9n-bux1ecb-dux1eef-liux1ec7u}}

Sau khi thu thập dữ liệu, trong thư mục dữ liệu sẽ có cấu trúc là 1 thư
mục dữ liệu duy nhất, trong đó tất cả là dữ liệu người dùng được định
dạng theo file json. đó.

\begin{figure}
\centering
\includegraphics{image/p3.png}
\caption{WorkFlow Hậu thu thập dữ liệu}
\end{figure}

Bước tiếp theo là tiến hành gộp chúng lại để hình thành bộ dữ liệu hợp
nhất ở mức người dùng. Và từ các trường dữ liệu sẵn có thì thực hiện tạo
thêm các trường dữ liệu mới dựa vào dữ liệu đã có như:

Quá trình này gồm các bước chính:

\begin{enumerate}
\def\labelenumi{\arabic{enumi}.}
\tightlist
\item
  \textbf{Xác định khóa chính (Primary Key)}

  \begin{itemize}
  \tightlist
  \item
    Sử dụng \texttt{user\_id} (hay \texttt{steamID}) làm định danh duy
    nhất cho mỗi người dùng.\\
  \item
    Mọi bảng dữ liệu phụ (achievements, games, groups, friends) đều được
    ánh xạ về bảng chính thông qua khóa này.
  \end{itemize}
\item
  \textbf{Chuẩn hóa dữ liệu trước khi gộp}

  \begin{itemize}
  \tightlist
  \item
    Đảm bảo các trường có định dạng đồng nhất:

    \begin{itemize}
    \tightlist
    \item
      Thời gian (\texttt{timecreated}) → chuyển sang dạng
      \textbf{datetime}
    \item
      Các cột nhị phân (\texttt{is\_old\_account},
      \texttt{is\_VN},\texttt{in\_game\_group},
      \texttt{in\_fan\_group},\texttt{in\_platform\_group}) → chuẩn về
      0/1.\\
    \item
      Các biến dạng chuỗi nhiều giá trị (như danh sách group) → chuyển
      sang vector đa nhãn
    \end{itemize}
  \end{itemize}
\item
  \textbf{Xóa các cột dữ liệu không cần thiết cho phân tích }

  \begin{itemize}
  \tightlist
  \item
    Xóa các cột như: ``realname'', ``profilestate'', ``personastate'',
    ``primaryclanid'', ``game\_achievements\_summary'',
    ``total\_playtime\_forever'', ``avatarfull''
  \end{itemize}
\item
  \textbf{Tạo bảng hợp nhất cuối cùng}

  \begin{itemize}
  \tightlist
  \item
    Mỗi hàng = 1 người dùng.\\
  \item
    Các cột đặc trưng:

    \begin{itemize}
    \tightlist
    \item
      \textbf{Thông tin cơ bản}: \texttt{user\_id},
      \texttt{account\_age\_years}, \texttt{loccountrycode}.\\
    \item
      \textbf{Hoạt động chơi game}: \texttt{total\_games},
      \texttt{total\_playtime\_hours},
      \texttt{avg\_playtime\_per\_game}.\\
    \item
      \textbf{Thành tựu}: \texttt{total\_achievements}.\\
    \item
      \textbf{Xã hội}: \texttt{friends\_count}, \texttt{group\_count},
      \texttt{social\_factor}.\\
    \item
      \textbf{Group Membership}: vector nhị phân từ \texttt{gid}.
    \end{itemize}
  \end{itemize}
\item
  \textbf{Tạo thêm các trường dữ liệu khác dựa vào dữ liệu đã có}
\end{enumerate}

\begin{itemize}
\item
  cột \texttt{timecreated} tạo ra được cọt
  \texttt{account\_age\_years},\texttt{account\_age\_da\ \ \ \ total\_playtime\_hours},
  \texttt{is\_old\_account}
\item
  Cột \texttt{loccountrycode} tạo ra được cột \texttt{is\_VN}, cột
  \texttt{locstatecode} tạo ra được cột \texttt{region} và cột
  \texttt{ten\_tinh}
\item
  Cột \texttt{game} tạo được cột \texttt{total\_game} ,
  \texttt{avg\_playtime\_per\_game}
\item
  Cột \texttt{group} tạo được cột
  \texttt{filtered\_all\_groups\_gid\_30} , \texttt{big\_group\_member}
  , \texttt{group\_count}, \texttt{in\_game\_group} ,
  \texttt{in\_fan\_group} ,\texttt{in\_platform\_group} ,
  \texttt{group\_name} , \texttt{social\_factor}
\item
  Cột \texttt{Friend} tạo được cột \texttt{friend\_steamids\_list},
  \texttt{friend\_count}
\item
  Cột \texttt{Achivement} tạo được cột \texttt{total\_achievements}
\end{itemize}

\begin{enumerate}
\def\labelenumi{\arabic{enumi}.}
\setcounter{enumi}{6}
\tightlist
\item
  \textbf{Bộ dữ liệu hoàn chỉnh}
\end{enumerate}

\begin{Shaded}
\begin{Highlighting}[]
\CommentTok{\# Đọc dữ liệu CSV}
\NormalTok{df }\OtherTok{\textless{}{-}} \FunctionTok{read.csv}\NormalTok{(}\StringTok{"data\_users.csv"}\NormalTok{)}

\CommentTok{\# Hiển thị tất cả tên cột}
\FunctionTok{colnames}\NormalTok{(df)}
\end{Highlighting}
\end{Shaded}

\begin{verbatim}
##  [1] "steamid"                    "user_id"                   
##  [3] "personaname"                "new_group_count"           
##  [5] "filtered_all_groups_gid_30" "group_names"               
##  [7] "timecreated"                "account_age_days"          
##  [9] "loccountrycode"             "locstatecode"              
## [11] "ten_tinh"                   "total_achievements"        
## [13] "total_games"                "total_playtime_hours"      
## [15] "big_group_member"           "avg_playtime_per_game"     
## [17] "account_age_years"          "is_old_account"            
## [19] "is_VN"                      "region"                    
## [21] "group_count"                "social_factor"             
## [23] "in_game_group"              "in_platform_group"         
## [25] "in_fan_group"               "friend_steamids_list"      
## [27] "friend_count"
\end{verbatim}

\hypertarget{iii.-tiux1ec1n-xux1eed-luxfd-vuxe0-tux1ea1o-ux111ux1eb7c-trux1b0ng-feature-engineering-preprocessing}{%
\section{III. Tiền xử lý và Tạo Đặc trưng (Feature Engineering \&
Preprocessing)}\label{iii.-tiux1ec1n-xux1eed-luxfd-vuxe0-tux1ea1o-ux111ux1eb7c-trux1b0ng-feature-engineering-preprocessing}}

Sau khi dữ liệu thô từ nhiều nguồn đã được gộp lại thành một
\textbf{bảng hợp nhất} ở mức người dùng, bước tiếp theo là \textbf{tiền
xử lý dữ liệu} và \textbf{xây dựng đặc trưng}. Đây là giai đoạn then
chốt, quyết định chất lượng của phân tích và mô hình sau này.

\hypertarget{tux1ea1o-ux111ux1eb7c-trux1b0ng-cuxf3-giuxe1-trux1ecb-phuxe2n-tuxedch}{%
\subsection{3.1 Tạo Đặc trưng có giá trị phân
tích}\label{tux1ea1o-ux111ux1eb7c-trux1b0ng-cuxf3-giuxe1-trux1ecb-phuxe2n-tuxedch}}

Một số đặc trưng quan trọng được tính toán từ dữ liệu gốc:

\begin{itemize}
\tightlist
\item
  \textbf{Tuổi tài khoản}

  \begin{itemize}
  \tightlist
  \item
    \texttt{account\_age\_years} = (ngày hiện tại --
    \texttt{timecreated}) / 365.\\
  \item
    Dùng để phân loại người dùng mới và cũ (\texttt{is\_old\_account}).
  \end{itemize}
\item
  \textbf{Chỉ số hoạt động chơi game}

  \begin{itemize}
  \tightlist
  \item
    \texttt{total\_games}: số lượng game sở hữu.\\
  \item
    \texttt{total\_playtime\_hours}: tổng thời gian chơi (giờ).\\
  \item
    \texttt{avg\_playtime\_per\_game}: trung bình giờ chơi trên mỗi
    game.
  \end{itemize}
\item
  \textbf{Chỉ số thành tích}

  \begin{itemize}
  \tightlist
  \item
    \texttt{total\_achievements}: tổng số achievement đã mở khóa.\\
  \item
    Phản ánh mức độ ``try-hard'' của người chơi.
  \end{itemize}
\item
  \textbf{Đặc trưng xã hội}

  \begin{itemize}
  \tightlist
  \item
    \texttt{friends\_count}: số lượng bạn bè
  \item
    \texttt{friend\_steamids\_list}: danh sách bạn bè.\\
  \item
    \texttt{group\_count}: số nhóm tham gia.\\
  \item
    \texttt{social\_factor}: kết hợp số bạn bè và số nhóm, nhằm đo lường
    mức độ gắn kết xã hội.
  \end{itemize}
\item
  \textbf{Đặc trưng nhị phân (One-hot/Binary Features)}

  \begin{itemize}
  \tightlist
  \item
    \texttt{is\_VN}: người dùng có mã quốc gia = Việt Nam.\\
  \item
    \texttt{in\_game\_group}, \texttt{in\_platform\_group},
    \texttt{in\_fan\_group}: xác định kiểu nhóm mà user tham gia.
  \end{itemize}
\end{itemize}

\hypertarget{xux1eed-luxfd-dux1eef-liux1ec7u-vuxe0-chuux1ea9n-huxf3a}{%
\subsection{3.2 Xử lý dữ liệu và Chuẩn
hóa}\label{xux1eed-luxfd-dux1eef-liux1ec7u-vuxe0-chuux1ea9n-huxf3a}}

Để đảm bảo dữ liệu ổn định cho mô hình:

\begin{enumerate}
\def\labelenumi{\arabic{enumi}.}
\tightlist
\item
  \textbf{Xử lý dữ liệu thiếu (Missing Values)}\\
  A. \textbf{Điền bằng Giá trị Trung vị (Median Imputation)} Sử dụng giá
  trị trung vị để điền vào các cột có tính chất số liên tục nhằm giảm
  thiểu ảnh hưởng của các giá trị ngoại lai (outliers):
\end{enumerate}

\begin{itemize}
\item
  Tổng số game (\texttt{total\_games}): Giá trị thiếu được điền bằng
  trung vị của toàn bộ cột.
\item
  Ngày tạo tài khoản (\texttt{timecreated}): Sau khi cột này được chuyển
  đổi sang định dạng datetime, các giá trị còn thiếu cũng được điền bằng
  trung vị của thời gian tạo tài khoản.
\end{itemize}

B. \textbf{Điền bằng Giá trị Số Mặc định (0 Imputation)}

\begin{itemize}
\item
  Các cột liên quan đến trạng thái hoặc số lượng vật lý được điền bằng
  0, ngụ ý rằng thông tin hoặc hoạt động đó không tồn tại:
\item
  Tổng thời gian chơi (\texttt{total\_playtime\_hours}): Nếu thiếu, gán
  bằng 0, giả định người chơi chưa có thời gian chơi được ghi nhận.
\item
  Tổng thành tích (\texttt{total\_achievements}): Nếu user không có
  thành tích, gán 0.
\item
  Trạng thái hồ sơ (\texttt{profilestate}) và Tổng số nhóm
  (\texttt{total\_groups}): Các cột này cũng được điền bằng 0 (giá trị
  mặc định hoặc không hoạt động).
\end{itemize}

C. \textbf{Điền bằng Chuỗi ký tự (String Imputation)}

Chúng tôi sử dụng các chuỗi ký tự mô tả để điền vào các giá trị bị thiếu
trong dữ liệu phân loại hoặc dữ liệu văn bản:

\begin{itemize}
\item
  Tên người dùng (\texttt{personaname}): Điền bằng `Unknown'.
\item
  Tên tỉnh/thành (\texttt{ten\_tinh}): Điền `UNKNOWN'.
\item
  Mã quốc gia (\texttt{loccountrycode}): Nếu thiếu, giá trị mặc định là
  `VN' (Việt Nam) hoặc `Unknown' trong các trường hợp khác.
\item
  Tóm tắt thành tích game (\texttt{game\_achievements\_summary}): Thay
  giá trị thiếu bằng None.
\item
  Danh sách Group GID (\texttt{all\_groups\_gid}): Điền chuỗi rỗng
  `{[}{]}' để chuẩn bị cho việc xử lý cấu trúc dữ liệu mảng sau này.
\item
  Xử lý Mã tỉnh/thành (\texttt{locstatecode}): Cột này có quy tắc riêng
  biệt. Nếu quốc gia là `VN' nhưng mã tỉnh/thành bị thiếu, chúng tôi gán
  một số ngẫu nhiên từ 1 đến 63 để mô phỏng 63 tỉnh/thành của Việt Nam.
  Với các quốc gia khác, giá trị được điền là `Unknown'.
\end{itemize}

\begin{enumerate}
\def\labelenumi{\arabic{enumi}.}
\setcounter{enumi}{1}
\tightlist
\item
  \textbf{Xử lý ngoại lệ (Outliers)}

  \begin{itemize}
  \tightlist
  \item
    Áp dụng Winsorization (giới hạn theo IQR) cho các biến có phân phối
    lệch mạnh (\texttt{total\_playtime\_hours},
    \texttt{total\_achievements},\texttt{total\_games},
    \texttt{account\_age\_years}, \texttt{avg\_playtime\_per\_game}).\\
  \item
    Giúp giảm ảnh hưởng của các user có hành vi bất thường (ví dụ: chơi
    50,000 giờ).\\
    \includegraphics{image/p8.png}
  \end{itemize}
\item
  \textbf{Chuẩn hóa dữ liệu số}

  \begin{itemize}
  \tightlist
  \item
    Các biến số (\texttt{total\_games},
    \texttt{avg\_playtime\_per\_game},\texttt{total\_playtime\_hours},
    \texttt{friend\_count}, \texttt{group\_count},
    \texttt{new\_group\_count}, \texttt{social\_factor},
    \texttt{account\_age\_days}, \texttt{account\_age\_years} \ldots)
    được chuẩn hóa bằng \textbf{StandardScaler} (Z-score).\\
  \item
    Đảm bảo các đặc trưng có cùng thang đo trước khi đưa vào mô hình
    (K-Means, PCA, \ldots).
  \end{itemize}
\end{enumerate}

\hypertarget{kux1ebft-quux1ea3-sau-tiux1ec1n-xux1eed-luxfd}{%
\subsection{3.3 Kết quả sau tiền xử
lý}\label{kux1ebft-quux1ea3-sau-tiux1ec1n-xux1eed-luxfd}}

Sau khi hoàn tất tiền xử lý, dữ liệu người dùng đã được làm sạch, chuẩn
hóa và sẵn sàng cho phân tích. Tất cả giá trị thiếu được xử lý bằng 0,
trung vị hoặc chuỗi mặc định; các ngoại lệ được giới hạn bằng phương
pháp Winsorization (IQR clipping) để loại bỏ ảnh hưởng của người dùng
bất thường. Toàn bộ các biến số được chuẩn hóa bằng StandardScaler, giúp
các đặc trưng có cùng thang đo. Kết quả thu được là một bảng dữ liệu hợp
nhất gồm khoảng 2.000 người dùng và hơn 25 đặc trưng phản ánh hành vi,
hoạt động và mức độ gắn kết xã hội. Dữ liệu sau xử lý đạt 0\% giá trị
thiếu, không còn ngoại lệ, phân phối cân bằng và sẵn sàng cho bước phân
tích giảm chiều (PCA) và phân cụm người dùng.

\hypertarget{trux1b0c-quan-hoa-dux1b0-liuxeau}{%
\subsection{3.4 TRực quan hóa dữ
liệu}\label{trux1b0c-quan-hoa-dux1b0-liuxeau}}

Sau khi tiền xử lý, dữ liệu được trực quan hóa nhằm đánh giá phân phối,
sự cân bằng và mối quan hệ giữa các đặc trưng:

\hypertarget{phuxe2n-phux1ed1i-cuxe1c-biux1ebfn-sux1ed1}{%
\subsubsection{3.4.1 Phân phối các biến
số:}\label{phuxe2n-phux1ed1i-cuxe1c-biux1ebfn-sux1ed1}}

Các đặc trưng định lượng như \texttt{total\_games},
\texttt{total\_playtime\_hours}, \texttt{avg\_playtime\_per\_game},
\texttt{friend\_count}, \texttt{group\_count}\ldots{} có phân phối lệch
phải (skewed), cho thấy phần lớn người dùng có mức độ hoạt động thấp,
chỉ một số ít có giá trị rất cao. Biến account\_age\_days có dạng gần
chuẩn, phản ánh sự phân bố đồng đều giữa các tài khoản mới và cũ.
\includegraphics{image/p12.png}

\hypertarget{phuxe2n-phux1ed1i-cuxe1c-biux1ebfn-phuxe2n-loux1ea1i}{%
\subsubsection{3.4.2 Phân phối các biến phân
loại:}\label{phuxe2n-phux1ed1i-cuxe1c-biux1ebfn-phuxe2n-loux1ea1i}}

Dữ liệu nhị phân như \texttt{is\_old\_account},
\texttt{in\_game\_group}, \texttt{in\_platform\_group},
\texttt{in\_fan\_group}\ldots{} nhìn chung khá cân bằng, ngoại trừ một
vài biến có xu hướng lệch (ví dụ, phần lớn người dùng thuộc nhóm
in\_fan\_group = 1). Biến region phân bố đều giữa các vùng, đảm bảo dữ
liệu không bị thiên lệch địa lý. \includegraphics{image/p10.png}

\hypertarget{tux1b0ux1a1ng-quan-giux1eefa-cuxe1c-biux1ebfn-sux1ed1}{%
\subsubsection{3.4.3 Tương quan giữa các biến
số:}\label{tux1b0ux1a1ng-quan-giux1eefa-cuxe1c-biux1ebfn-sux1ed1}}

Ma trận tương quan (heatmap) cho thấy các đặc trưng
\texttt{group\_count}, \texttt{new\_group\_count} và
\texttt{social\_factor} có mức tương quan cao (≈ 0.9), phản ánh chúng
cùng mô tả mức độ tham gia nhóm xã hội. Trong khi đó, các biến như
\texttt{total\_games}, \texttt{total\_playtime\_hours},
\texttt{friend\_count}, \texttt{account\_age\_days} có tương quan yếu,
thể hiện tính độc lập, phù hợp cho các mô hình phân tích đa chiều.
\includegraphics{image/p11.png}

\hypertarget{liuxean-kuxeat-xa-huxf4i}{%
\subsection{3.4 Liên Kết Xã Hội}\label{liuxean-kuxeat-xa-huxf4i}}

Vì Steam là một nền tảng chơi game online đa nề tảng, gióng như một mạng
xã hội của các gamer nên người dùng không chỉ thể hiện hoạt động cá nhân
thông qua việc chơi game mà còn hình thành mối quan hệ xã hội với những
người chơi khác thông qua các nhóm, bạn bè, hoặc cộng đồng chung. Do đó,
việc phân tích liên kết xã hội giúp hiểu sâu hơn về cấu trúc tương tác,
mức độ kết nối, và ảnh hưởng xã hội của từng người dùng trong hệ sinh
thái này.

\hypertarget{suxf4-node-va-suxf4-canh}{%
\subsubsection{3.4.1 Số node và số
cạnh}\label{suxf4-node-va-suxf4-canh}}

Sau khi xử lý dữ liệu, số lượng node trong mạng xã hội tương ứng với số
người dùng hợp lệ (không trùng lặp), còn số lượng cạnh được tính dựa
trên tổng số quan hệ bạn bè có thật giữa các người dùng có trong tập dữ
liệu.

Trong bộ dữ liệu này: số node là 1482 và số cạnh là: 57796

\hypertarget{ux111uxf4-thi-biuxeau-diuxean-lien-kuxeat-xa-huxf4i}{%
\subsubsection{3.4.2 Đồ thị biểu diễn lien kết xã
hội}\label{ux111uxf4-thi-biuxeau-diuxean-lien-kuxeat-xa-huxf4i}}

\begin{itemize}
\tightlist
\item
  Để trực quan hóa cấu trúc kết nối, đồ thị mạng được vẽ với mỗi node là
  một người dùng, và cạnh thể hiện mối quan hệ bạn bè.
\end{itemize}

\hypertarget{a-ux111uxf4-thi-biuxeau-diuxean-sux1b0-lien-kuxeat-cua-100-users-ux111uxe2u-tiuxean}{%
\paragraph{\texorpdfstring{\textbf{A đồ thị biểu diễn sứ lien kết của
100 users đầu
tiên}}{A đồ thị biểu diễn sứ lien kết của 100 users đầu tiên}}\label{a-ux111uxf4-thi-biuxeau-diuxean-sux1b0-lien-kuxeat-cua-100-users-ux111uxe2u-tiuxean}}

\includegraphics{image/p5.png}

Phần lớn các nút (người dùng) tập trung ở vùng trung tâm, cho thấy có
một cộng đồng chính trong đó nhiều người có mối liên kết lẫn nhau.

Một số nút nằm rải rác ở ngoài, thể hiện các người dùng cô lập hoặc chỉ
có ít kết nối (thậm chí không có bạn chung).

Mạng lưới chưa quá dày đặc, tức là mức độ kết nối giữa người dùng vẫn
còn hạn chế --- nhiều người chỉ có 1--2 liên kết.

Sự xuất hiện của một cụm lớn trung tâm cho thấy có khả năng tồn tại một
nhóm người dùng hoạt động mạnh hoặc có sở thích chung.

Một vài thành phần tách biệt (disconnected components) ở rìa biểu đồ thể
hiện những nhóm nhỏ hoặc cá nhân không liên hệ với mạng chính.

\hypertarget{b-ux111uxf4-thi-biuxeau-diuxean-ux111uxe2y-ux111u-sux1b0-lien-kuxeat-giux1b0a-cac-ngux1b0ux1a1i-chux1a1i}{%
\paragraph{\texorpdfstring{\textbf{B đồ thị biểu diễn đầy đủ sự lien kết
giữa các người
chơi:}}{B đồ thị biểu diễn đầy đủ sự lien kết giữa các người chơi:}}\label{b-ux111uxf4-thi-biuxeau-diuxean-ux111uxe2y-ux111u-sux1b0-lien-kuxeat-giux1b0a-cac-ngux1b0ux1a1i-chux1a1i}}

\includegraphics{image/p6.png} Mạng lưới có mật độ kết nối cao, thể hiện
qua vùng trung tâm dày đặc, nơi nhiều người dùng có mối quan hệ chồng
chéo.

Xuất hiện một cộng đồng chính rất lớn, trong đó hầu hết người dùng đều
liên kết với nhau, cho thấy tính liên thông mạnh giữa các tài khoản.

Một số nút nằm ở rìa hoặc tách biệt, biểu hiện cho các người dùng ít bạn
hoặc chỉ kết nối với một số thành viên trong mạng.

Cấu trúc này gợi ý rằng hệ thống có một mạng xã hội trung tâm (core
network), với các nhóm nhỏ hơn lan tỏa xung quanh.

\hypertarget{c-la-biuxeau-ux111uxf4-tux1b0ux1a1ng-quan-thuxea-hiuxean-sux1b0-tux1b0ux1a1ng-ux111uxf4ng-cua-ngux1b0ux1a1i-dung-thuxea-hiuxean-qua-cac-biuxean-dux1b0-liuxeau}{%
\paragraph{\texorpdfstring{\textbf{C là biểu đồ tương quan thể hiện sự
tương đồng của người dùng thể hiện qua các biến dữ
liệu}}{C là biểu đồ tương quan thể hiện sự tương đồng của người dùng thể hiện qua các biến dữ liệu}}\label{c-la-biuxeau-ux111uxf4-tux1b0ux1a1ng-quan-thuxea-hiuxean-sux1b0-tux1b0ux1a1ng-ux111uxf4ng-cua-ngux1b0ux1a1i-dung-thuxea-hiuxean-qua-cac-biuxean-dux1b0-liuxeau}}

Vì nếu vẽ ra toàn bộ 1498 người thì rõ ràng biểu đồ rất khó nhìn, nên ở
đây chỉ biểu diễn trên 10 người đầu tiên.

Các yếu tố để xét sự tương đồng là: ``loccountrycode'', ``region'',
``ten\_tinh'',``in\_game\_group'', ``in\_platform\_group'',
``in\_fan\_group'', ``friend\_steamids\_list''. bao gồm: Quốc gia, khu
vực (Bắc/Trung/Nam), tên tỉnh/thành phố, nằm trong các nhóm đã được phân
loại, và có bạn chung Cách tính điểm trên biểu đồ tương quan là: Nếu 2
người có 1 bạn chung thì cộng 0.2 điểm, các yếu tốc khác chung sẽ đươc
cộng 1 điểm. \includegraphics{image/p7.png}

Biểu đồ heatmap thể hiện điểm tương đồng giữa 10 người dùng đầu tiên
trong tập dữ liệu. Quan sát cho thấy:

Các ô màu đậm (xanh đậm) biểu thị cặp người dùng có mức độ tương đồng
cao, tức là họ có nhiều đặc điểm giống nhau (ví dụ: cùng quốc gia, cùng
vùng, tham gia nhóm giống nhau hoặc có bạn chung).

Các ô màu nhạt (xanh nhạt hoặc vàng nhạt) biểu thị mức độ tương đồng
thấp, tức là hai người dùng có ít điểm chung hơn.

Đường chéo chính có giá trị 0 vì mỗi người không so sánh với chính mình.

Một số nhóm người dùng (ví dụ User3--User4--User6--User9) có điểm tương
đồng cao lặp lại, cho thấy họ có thể thuộc cùng cộng đồng hoặc nhóm hành
vi tương tự.

Ngược lại, các người dùng như User7 có giá trị tương đồng thấp hơn với
hầu hết người khác, cho thấy họ ít điểm chung hơn với phần còn lại.

\hypertarget{iv.phuxe2n-cux1ee5m-cux1ed9ng-ux111ux1ed3ng.}{%
\subsection{IV.Phân cụm cộng
đồng.}\label{iv.phuxe2n-cux1ee5m-cux1ed9ng-ux111ux1ed3ng.}}

Mục tiêu chính là xác định và phân loại các nhóm người dùng có hành vi
và đặc điểm tương đồng dựa trên các thông tin về hoạt động, tương tác xã
hội và sở thích cá nhân. Việc này không chỉ giúp hiểu rõ cấu trúc và đặc
điểm của cộng đồng người dùng mà còn tạo tiền đề quan trọng cho các phân
tích nâng cao sau này. Bằng cách xác định các nhóm đồng nhất, chúng ta
có thể:

A. \textbf{PHÂN CỤM HOẠT ĐỘNG}

Chỉ tập trung vào các đặc trưng về hoạt động chơi game: số game chơi,
tổng thời gian chơi, thời gian trung bình cho mỗi game\ldots{}

Mục tiêu: phân loại người dùng dựa trên mức độ hoạt động trong game, ví
dụ casual vs hardcore.

Ứng dụng: gợi ý game mới, chương trình reward, hoặc đánh giá engagement.

B. \textbf{PHÂN CỤM XÃ HỘI}

Tập trung vào tương tác xã hội: số bạn bè, số nhóm tham gia,
social\_factor, thành viên nhóm lớn\ldots{}

Mục tiêu: phân loại người dùng dựa trên mức độ kết nối và tham gia cộng
đồng.

Ứng dụng: gợi ý kết nối bạn bè, event xã hội, phân tích mạng lưới cộng
đồng.

\hypertarget{biux1ec3u-diux1ec5n-dux1eef-liux1ec7u-ngux1b0ux1eddi-duxf9ng-nhuxf3m}{%
\subsubsection{4.1 Biểu diễn dữ liệu người dùng --
nhóm}\label{biux1ec3u-diux1ec5n-dux1eef-liux1ec7u-ngux1b0ux1eddi-duxf9ng-nhuxf3m}}

Trong cộng đồng người dùng, mỗi cá nhân có thể tham gia nhiều loại nhóm
khác nhau, chẳng hạn như nhóm chơi game, nhóm xã hội, nhóm fan hoặc các
nhóm sở thích chuyên biệt khác. Để mô hình hóa sự tham gia này, ta biểu
diễn dữ liệu dưới dạng ma trận Người dùng -- Nhóm:

Hàng (rows): mỗi hàng đại diện cho một người dùng duy nhất trong hệ
thống.

Cột (columns): mỗi cột tương ứng với một nhóm cụ thể trong cộng đồng.

Ô (i, j): giá trị của ô là 1 nếu người dùng thứ i thuộc nhóm thứ j, và 0
nếu không tham gia.

Cách biểu diễn này giúp mô hình nhận biết sự tương đồng giữa các người
dùng dựa trên việc chia sẻ nhóm chung. Những người dùng tham gia nhiều
nhóm giống nhau sẽ có các hàng ma trận gần nhau, phản ánh mối liên hệ về
sở thích, hành vi hoặc mức độ tham gia trong cộng đồng. Nhờ ma trận này,
các thuật toán phân cụm có thể phát hiện ra các nhóm người dùng đồng
nhất, phục vụ cho các phân tích nâng cao và các ứng dụng như gợi ý kết
nối, cá nhân hóa trải nghiệm hay đánh giá mức độ gắn kết trong cộng
đồng.

\hypertarget{giux1ea3m-chiux1ec1u-vuxe0-chuux1ea9n-huxf3a-dux1eef-liux1ec7u}{%
\subsubsection{4.2 Giảm chiều và chuẩn hóa dữ
liệu}\label{giux1ea3m-chiux1ec1u-vuxe0-chuux1ea9n-huxf3a-dux1eef-liux1ec7u}}

\begin{itemize}
\tightlist
\item
  Do số nhóm rất lớn và dữ liệu thưa, cần giảm chiều để tăng hiệu quả
  phân tích.\\
\item
  Phương pháp \textbf{PCA (Phân tích thành phần chính)} giúp nén thông
  tin nhưng vẫn giữ lại sự khác biệt quan trọng.\\
\item
  Sau đó chuẩn hóa dữ liệu để các đặc trưng hành vi (số giờ chơi, số
  game, thành tích\ldots) có cùng thang đo.
\end{itemize}

\hypertarget{phuxe2n-cum-bux103ng-k-means-va-dbscan}{%
\subsubsection{4.3 Phân cụm bằng K-means và
DBscan}\label{phuxe2n-cum-bux103ng-k-means-va-dbscan}}

Để phân loại người dùng thành các nhóm có hành vi và đặc điểm tương
đồng, ta áp dụng hai thuật toán phân cụm phổ biến:

\textbf{K-Means Clustering}

Chia người dùng thành K cụm dựa trên khoảng cách giữa các đặc trưng.

Số cụm K được lựa chọn dựa trên Elbow method, tức quan sát điểm gãy
trong đồ thị tổng bình phương khoảng cách trong cụm (Inertia) để xác
định số cụm hợp lý.

Phương pháp này giúp nhóm những người dùng có hành vi tương tự nhau
thành các cụm đồng nhất.

\textbf{DBSCAN Clustering}

Phân cụm dựa trên mật độ dữ liệu, tự động xác định các cụm có hình dạng
khác nhau và phát hiện các điểm dữ liệu nhiễu (outlier).

Thuật toán này đặc biệt hữu ích khi các cụm không có kích thước hoặc
hình dạng đều nhau.

Sử dụng kết hợp cả hai thuật toán cho phép vừa phân cụm theo cấu trúc
hình học (K-Means), vừa phát hiện các nhóm đặc biệt và điểm ngoại lệ
trong dữ liệu (DBSCAN).

\hypertarget{thux1b0c-hiuxean-phuxe2n-cum}{%
\subsubsection{4.4 Thực hiện phân
cụm}\label{thux1b0c-hiuxean-phuxe2n-cum}}

\hypertarget{xuxe1c-ux111ux1ecbnh-sux1ed1-cux1ee5m-bux1eb1ng-elbow-method}{%
\paragraph{4.4.1 Xác định số cụm bằng Elbow
Method}\label{xuxe1c-ux111ux1ecbnh-sux1ed1-cux1ee5m-bux1eb1ng-elbow-method}}

Để chọn số cụm hợp lý cho K-Means, ta áp dụng Elbow method:

Vẽ đồ thị Inertia (tổng bình phương khoảng cách trong cụm) theo số cụm
K.

Quan sát điểm gãy (elbow) trong đồ thị, tại đó việc tăng số cụm không
còn giảm đáng kể Inertia nữa.

Điểm gãy này được chọn làm số cụm tối ưu, giúp K-Means phân loại dữ liệu
hiệu quả mà không quá nhiều hoặc quá ít cụm.
\includegraphics{image/p13.png}

\textbf{Elbow - Social (Phân Cụm Xã Hội)} Biểu đồ này thể hiện Inertia
giảm dần khi số cụm (k) tăng lên.

Điểm Uốn Khả Thi: Điểm uốn (elbow) rõ ràng nhất có thể được nhìn thấy
tại k=3 hoặc k=4.

Sự giảm Inertia từ k=2 xuống k=3 là rất lớn.

Sự giảm từ k=3 xuống k=4 vẫn đáng kể.

Từ k=4 trở đi, đường cong trở nên thoai thoải hơn, cho thấy lợi ích biên
của việc thêm cụm giảm đi.

Lựa Chọn Khuyến Nghị: Lựa chọn tối ưu nhất có xu hướng là k=3 để đạt
được sự cân bằng giữa độ nén của cụm.

\textbf{Elbow - Activity (Phân Cụm Hoạt Động)}

Biểu đồ này cũng cho thấy Inertia giảm nhanh chóng ở k nhỏ.

Điểm Uốn Khả Thi: Điểm uốn (elbow) rõ rệt có vẻ xuất hiện tại k=3 hoặc
k=4.

Sự giảm Inertia từ k=2 xuống k=3 là rất mạnh.

Sự giảm từ k=3 xuống k=4 vẫn khá dốc.

Tương tự như biểu đồ Social, chọn k=3,sẽ đạt được độ ổn định hơn.

\hypertarget{kuxeat-qua-phuxe2n-cum}{%
\paragraph{4.4.2 Kết quả phân cụm}\label{kuxeat-qua-phuxe2n-cum}}

\includegraphics{image/p14.png} \includegraphics{image/p15.png}
\includegraphics{image/p16.png} \includegraphics{image/p17.png}

\hypertarget{so-sanh-hai-muxf4-hinh}{%
\subsubsection{4.4 So sánh hai mô hình}\label{so-sanh-hai-muxf4-hinh}}

\begin{longtable}[]{@{}
  >{\raggedright\arraybackslash}p{(\columnwidth - 4\tabcolsep) * \real{0.1127}}
  >{\raggedright\arraybackslash}p{(\columnwidth - 4\tabcolsep) * \real{0.4507}}
  >{\raggedright\arraybackslash}p{(\columnwidth - 4\tabcolsep) * \real{0.4366}}@{}}
\toprule()
\begin{minipage}[b]{\linewidth}\raggedright
Tiêu chí
\end{minipage} & \begin{minipage}[b]{\linewidth}\raggedright
K-Means
\end{minipage} & \begin{minipage}[b]{\linewidth}\raggedright
DBSCAN
\end{minipage} \\
\midrule()
\endhead
Nguyên tắc hoạt động & Chia dữ liệu thành k cụm dựa trên khoảng cách
Euclide đến tâm cụm , giả định cụm có hình dạng cầu. & Phân nhóm dữ liệu
dựa trên mật độ , tìm kiếm các khu vực dày đặc bị tách biệt bởi các khu
vực có mật độ thấp. \\
Xử lý Nhiễu & Gán mọi điểm vào một cụm bất kỳ, ngay cả điểm ngoại lai. &
Xác định rõ ràng các điểm có mật độ thấp là Nhiễu (-1). \\
Hiệu suất trên Dữ liệu này & Tốt hơn. Tạo ra các cụm có kích thước tương
đối cân bằng, dễ diễn giải, và phù hợp với mục tiêu tạo phân khúc thị
trường . & Kém hiệu quả. Thường coi phần lớn dữ liệu là một cụm duy nhất
(Cụm 0) hoặc là Nhiễu (-1), đặc biệt trong Phân cụm Hoạt động và Xã
hội. \\
Yêu cầu Tham số & Chỉ cần xác định k (số cụm), được hỗ trợ bởi phương
pháp Elbow. & Cần tinh chỉnh ϵ (bán kính) và MinSamples (mật độ tối
thiểu), rất nhạy cảm với sự thay đổi mật độ dữ liệu. \\
\bottomrule()
\end{longtable}

\textbf{Phân Tích Chuyên Sâu theo Kết Quả}

\textbf{1. Ưu thế của K-Means} Tính Ứng dụng Cao: K-Means đã thành công
trong việc tạo ra các phân khúc có ý nghĩa trong cả hai hạng mục (Xã
hội, Hoạt động).

Trong Phân cụm, K-Means đã phân tách rõ ràng Cụm 3 (màu đỏ), một nhóm
người dùng có hồ sơ cực kỳ khác biệt (PC2 rất thấp), khỏi khối dữ liệu
chính. Điều này là cực kỳ quan trọng cho việc mô tả hồ sơ.

Độ Ổn định và Tính Diễn giải: Do buộc phải chia dữ liệu thành k cụm,
K-Means cung cấp cái nhìn tổng thể về cách cộng đồng được phân bổ, ngay
cả khi các cụm có sự chồng chéo.

\textbf{2. Hạn chế của DBSCAN} Quá Nhạy Cảm với Mật Độ: Các biểu đồ cho
thấy dữ liệu người dùng có mật độ rất dày và đồng nhất (kết quả của việc
chuẩn hóa).

Trong Phân cụm Hoạt động, DBSCAN gần như gán toàn bộ dữ liệu vào một cụm
duy nhất (Cụm 0), cho thấy nó không thể tìm thấy các ``khoảng trống'' để
tách người dùng thành các cấp độ hoạt động khác nhau.

Trong Phân cụm Xã hội, DBSCAN lại quá nhạy cảm và coi phần lớn là nhiễu
(-1).

Không Phù hợp với Mục tiêu: Mục tiêu là tạo ra các phân khúc có thể định
nghĩa được để làm đầu vào cho mô hình dự đoán. DBSCAN, bằng cách tạo ra
quá nhiều nhiễu hoặc quá ít cụm có ý nghĩa, không đáp ứng được mục tiêu
này.

\hypertarget{uxfd-nghux129a-kux1ebft-quux1ea3-phuxe2n-cux1ee5m}{%
\subsubsection{4.5 Ý nghĩa kết quả phân
cụm}\label{uxfd-nghux129a-kux1ebft-quux1ea3-phuxe2n-cux1ee5m}}

\textbf{Cụm Xã Hội}

\begin{longtable}[]{@{}
  >{\raggedright\arraybackslash}p{(\columnwidth - 6\tabcolsep) * \real{0.0798}}
  >{\raggedright\arraybackslash}p{(\columnwidth - 6\tabcolsep) * \real{0.1681}}
  >{\raggedright\arraybackslash}p{(\columnwidth - 6\tabcolsep) * \real{0.4076}}
  >{\raggedright\arraybackslash}p{(\columnwidth - 6\tabcolsep) * \real{0.3445}}@{}}
\toprule()
\begin{minipage}[b]{\linewidth}\raggedright
Cụm (KMeans)
\end{minipage} & \begin{minipage}[b]{\linewidth}\raggedright
Phân Bố trên PCA 2D
\end{minipage} & \begin{minipage}[b]{\linewidth}\raggedright
Hồ Sơ Đặc Trưng (Cluster Profile)
\end{minipage} & \begin{minipage}[b]{\linewidth}\raggedright
Vai Trò Cộng Đồng
\end{minipage} \\
\midrule()
\endhead
Cụm 0 (Xanh Dương) & Trung tâm, rộng, chồng chéo với Cụm 1 & Người Dùng
Cân Bằng/Tiêu Chuẩn: Mạng lưới xã hội vừa phải và ổn định. Tương tác và
tham gia nhóm ở mức trung bình. & Nền tảng: Tạo sự ổn định và là cầu nối
cho các cụm khác. \\
Cụm 1 (Cam) & Phía trên \& phải (PC1 cao) & Người Xã Hội Năng Động : Số
lượng bạn bè và nhóm tham gia rất cao. Là những người dùng có xu hướng
kết nối mạnh mẽ, làm trung tâm giao tiếp. & Chủ chốt: Quan trọng để lan
truyền thông tin, kết nối mạng lưới, và thúc đẩy hoạt động xã hội. \\
Cụm 2 (Xanh Lá) & Phía trên \& trái (PC1 thấp) & Người Ít Giao Tiếp Xã
Hội : Số lượng bạn bè và nhóm tham gia thấp hoặc rất chọn lọc. Tập trung
vào game/nội dung hơn là mở rộng mạng lưới. & Người dùng Thầm Lặng: Cần
chiến lược can thiệp để thúc đẩy tương tác xã hội. \\
\bottomrule()
\end{longtable}

\textbf{Cụm Hoạt Động}

\begin{longtable}[]{@{}
  >{\raggedright\arraybackslash}p{(\columnwidth - 6\tabcolsep) * \real{0.0643}}
  >{\raggedright\arraybackslash}p{(\columnwidth - 6\tabcolsep) * \real{0.1786}}
  >{\raggedright\arraybackslash}p{(\columnwidth - 6\tabcolsep) * \real{0.3857}}
  >{\raggedright\arraybackslash}p{(\columnwidth - 6\tabcolsep) * \real{0.3714}}@{}}
\toprule()
\begin{minipage}[b]{\linewidth}\raggedright
Cụm
\end{minipage} & \begin{minipage}[b]{\linewidth}\raggedright
Vị Trí trên PCA 2D (Tương đối)
\end{minipage} & \begin{minipage}[b]{\linewidth}\raggedright
Hồ Sơ Đặc Trưng (Giả định)
\end{minipage} & \begin{minipage}[b]{\linewidth}\raggedright
Ý Nghĩa Cộng Đồng
\end{minipage} \\
\midrule()
\endhead
Cụm 0 (Xanh Dương) & Trung tâm và rải rác & Người Chơi Casual (Thấp): Có
tổng giờ chơi và số game thấp (có thể là người chơi ít thường xuyên hoặc
chỉ mới bắt đầu). & Nhóm có nguy cơ rời bỏ cao, cần đề xuất nội dung
hoặc nhóm chơi hấp dẫn. \\
Cụm 1 (Cam) & Phía trên bên phải (PC1 cao, PC2 cao/trung bình) & Người
Chơi Core/Hardcore (Cao): Có tổng giờ chơi và số game rất cao. Họ là
những người dành nhiều thời gian nhất cho nền tảng. & High-Value Users
(Người dùng giá trị cao), đóng góp lớn vào sự duy trì hoạt động của nền
tảng. \\
Cụm 2 (Xanh Lá) & Phía trên bên trái (PC1 thấp, PC2 cao) & Người Chơi
Trung Bình/Chuyên Sâu: Có mức độ hoạt động vừa phải; có thể là những
người chơi thường xuyên nhưng không nhiều giờ như Cụm 1, hoặc những
người chơi nhiều giờ nhưng chỉ tập trung vào một số game nhất định. &
Nhóm có tiềm năng phát triển, cần khuyến khích thử nghiệm game mới hoặc
tương tác cộng đồng. \\
\bottomrule()
\end{longtable}

\hypertarget{v.-kuxeat-luuxe2n}{%
\section{V. Kết luận}\label{v.-kuxeat-luuxe2n}}

Quy trình phân cụm đã thành công trong việc phân đoạn người dùng thành
các nhóm có ý nghĩa bằng thuật toán K-Means, vượt trội so với DBSCAN
trong việc tạo ra các phân khúc dễ diễn giải. Cụ thể, phân tích đã xác
định được các hồ sơ quan trọng thông qua các cụm được hỗ trợ bởi phương
pháp Elbow. Các nhãn cụm này đóng vai trò là biến đại diện mạnh mẽ, cô
đọng hành vi phức tạp của người dùng. Việc đưa các nhãn này vào mô hình
dự đoán sẽ đơn giản hóa tính toán tương đồng, xác định tính bổ sung và
cá nhân hóa gợi ý kết nối, từ đó nâng cao độ chính xác trong việc dự
đoán và thúc đẩy các mối quan hệ tích cực giữa những người chơi.

\end{document}
